\topic{本章实验、故意破坏与练习}

\begin{experimentbox}{记录一次真实抢占}
让两个 Thread 使用不同通用寄存器和 XMM0 模式，跨 100 次抢占后核对。每次
真正切换只记 from、to、tick 和原因。
\end{experimentbox}

\begin{faultinjectionbox}{忘记更新 TSS.RSP0}
切到目标 Thread 却保留旧内核栈顶。预期下一次 CPL3 中断把帧压到错误 Thread
的栈。恢复切换顺序后，帧地址应落在目标栈范围。
\end{faultinjectionbox}

\begin{pitfallbox}
Ready 与队列位置同步改变；系统调用入口先找到当前 CpuLocal；用户 RSP 不能
直接当内核栈使用。
\end{pitfallbox}

\textbf{练习}
\begin{enumerate}[leftmargin=2.2em]
  \item 时间片到但 Ready queue 空时是否切换？
  \item syscall 的用户返回 RIP 通常来自哪个寄存器？
  \item 为什么 CR3 要和 Thread 所属 Process 一起切换？
\end{enumerate}
\begin{answerbox}
1. 不切；2. RCX（按 SYSCALL/SYSRET 约定）；3. 相同虚拟地址在不同 Process
中可能映射不同物理页。
\end{answerbox}
